\documentclass[12pt]{article}
\usepackage{amsmath, amssymb, amsthm}
\usepackage{geometry}
\geometry{margin=1in}

\newtheorem{exercise}{Exercise}
\theoremstyle{remark}
\newtheorem*{solution}{Solution}
\newtheorem*{remark}{Remark}

\title{Random Walks --- Week 7 Piazza Exercises}
\author{CUNY, Spring 2026}
\date{}

\begin{document}
\maketitle

\section*{Week 7: Conditional Expectation}

Let $(X, Y)$ be uniformly distributed on the unit disk
$D = \{(x,y) \in \mathbb{R}^2 : x^2 + y^2 \le 1\}$.
The joint density is $f_{X,Y}(x,y) = 1/\pi$ on $D$.

\medskip\noindent
The marginal density of $Y$ is
\[
  f_Y(y) = \int_{-\sqrt{1-y^2}}^{\sqrt{1-y^2}} \frac{1}{\pi}\,dx
  = \frac{2\sqrt{1-y^2}}{\pi}, \qquad y \in [-1,1].
\]
Given $Y = y$, the conditional density of $X$ is
\[
  f_{X \mid Y}(x \mid y) = \frac{f_{X,Y}(x,y)}{f_Y(y)}
  = \frac{1/\pi}{2\sqrt{1-y^2}/\pi}
  = \frac{1}{2\sqrt{1-y^2}},
  \qquad x \in \bigl[-\sqrt{1-y^2},\,\sqrt{1-y^2}\bigr],
\]
so $X \mid Y = y$ is uniform on $[-\sqrt{1-y^2},\,\sqrt{1-y^2}]$.

% ----------------------------------------------------------------
\begin{exercise}[Conditional probability $P(X \le Y \mid Y)$]
Compute $P(X \le Y \mid Y)$ as an explicit function of $Y$.
\end{exercise}

\begin{solution}
Given $Y = y$ with $y \in [-1,1]$, $X$ is uniform on $[-\sqrt{1-y^2},\,\sqrt{1-y^2}]$.
Let $a(y) = -\sqrt{1-y^2}$ and $b(y) = \sqrt{1-y^2}$.

We compute $P(X \le y \mid Y = y)$ by integrating the uniform density:
\[
  P(X \le y \mid Y = y)
  = \frac{(y \wedge b(y)) - a(y)}{b(y) - a(y)}
  = \frac{(y \wedge \sqrt{1-y^2}) + \sqrt{1-y^2}}{2\sqrt{1-y^2}}.
\]
We consider three cases depending on the sign of $y - \sqrt{1-y^2}$:

\textbf{Case 1:} $y \in (-\sqrt{2}/2,\, \sqrt{2}/2)$ (so $|y| < \sqrt{1-y^2}$, i.e.\
$y^2 < 1/2$).  Then $y \in (a(y), b(y))$, so
\[
  P(X \le y \mid Y = y) = \frac{y + \sqrt{1-y^2}}{2\sqrt{1-y^2}}.
\]

\textbf{Case 2:} $y \in [\sqrt{2}/2,\, 1]$ (so $y \ge \sqrt{1-y^2}$).
Then $y \ge b(y)$, meaning $X \le b(y) \le y$ always, so
\[
  P(X \le y \mid Y = y) = 1.
\]

\textbf{Case 3:} $y \in [-1,\, -\sqrt{2}/2]$ (so $y \le -\sqrt{1-y^2} = a(y)$).
Then $y \le a(y)$, meaning $X \ge a(y) \ge y$ always, so
\[
  P(X \le y \mid Y = y) = 0.
\]

\medskip\noindent
Combining, as a function of $Y$:
\[
  \boxed{P(X \le Y \mid Y) = \begin{cases}
    0 & \text{if } Y \in [-1,\,-\sqrt{2}/2], \\[4pt]
    \dfrac{Y + \sqrt{1-Y^2}}{2\sqrt{1-Y^2}} & \text{if } Y \in (-\sqrt{2}/2,\,\sqrt{2}/2), \\[8pt]
    1 & \text{if } Y \in [\sqrt{2}/2,\,1].
  \end{cases}}
\]
Note: the boundary points have $P(Y = \pm\sqrt{2}/2) = 0$ so their assignment is
irrelevant (the conditional expectation is defined only a.s.).
\end{solution}

% ----------------------------------------------------------------
\begin{exercise}[Conditional expectation $E(X \mid Y)$]
Find $E(X \mid Y)$.
\end{exercise}

\begin{solution}
Given $Y = y$, $X$ is uniform on $[-\sqrt{1-y^2},\,\sqrt{1-y^2}]$, a symmetric interval
about $0$.  Therefore
\[
  E(X \mid Y = y) = \frac{-\sqrt{1-y^2} + \sqrt{1-y^2}}{2} = 0.
\]
Hence $E(X \mid Y) = 0$ a.s.

\medskip\noindent\textit{Verification:} We can check directly that $\int_A E(X \mid Y)\,dP = \int_A X\,dP$ for $A \in \sigma(Y)$.  Since $E(X \mid Y) = 0$ and $X$ is symmetric about $0$ given any $y$, both integrals equal $0$.
\end{solution}

% ----------------------------------------------------------------
\begin{exercise}[Conditional second moment $E(X^2 \mid Y)$]
Find $E(X^2 \mid Y)$.
\end{exercise}

\begin{solution}
Given $Y = y$, $X \sim \mathrm{Unif}[-\sqrt{1-y^2},\,\sqrt{1-y^2}]$.
For a uniform distribution on $[-c, c]$, the second moment is
\[
  E[X^2 \mid Y = y]
  = \int_{-c}^{c} x^2 \cdot \frac{1}{2c}\,dx
  = \frac{1}{2c} \cdot \frac{2c^3}{3}
  = \frac{c^2}{3},
\]
with $c = \sqrt{1 - y^2}$.  Therefore
\[
  \boxed{E(X^2 \mid Y) = \frac{1 - Y^2}{3}.}
\]
\textit{Verification:} $E\bigl[E(X^2 \mid Y)\bigr] = E\bigl[\frac{1-Y^2}{3}\bigr]
= \frac{1}{3}\bigl(1 - E[Y^2]\bigr)$.  Since $E[Y^2] = \frac{1}{\pi}\int_D y^2\,dx\,dy
= \frac{1}{4}$ (for the uniform distribution on the unit disk), we get
$E[X^2 \mid Y] = \frac{1}{3}(1 - \frac{1}{4}) = \frac{1}{4}$.
Directly, $E[X^2] = \frac{1}{4}$ by symmetry, which confirms the calculation.
\end{solution}

\end{document}
